\documentclass[12pt]{article}
\usepackage{amsmath,amssymb,geometry,setspace,url}
\geometry{margin=1in}
\setstretch{1.2}
\begin{document}
\title{SKANN-SSL --- DSP \& Sampling Standards (Stage 0 Preprocessing)}
\author{Oravont Systems LLP}
\date{}
\maketitle

SKANN Document C
DSP and Sampling Logic for SKANN Stage 0
(Final Structured Draft)

Oravont Systems LLP

\section*{Introduction}
Document C defines the digital signal processing (DSP) conventions used in the SKANN project. These conventions ensure that all synthetic and recorded underwater signals are:

comparable in sampling rate,

comparable in amplitude scale,

represented in a consistent spectral framework,

compatible with machine-learning front-ends,

correctly aligned with Documents A (acoustics) and B (ambient noise).

SKANN Stage 0 (Preprocessing \& Standardisation) uses the definitions in this document to:

resample incoming audio,

normalise amplitude,

compute FFT size and bin spacing,

segment data,

prepare tensors for the learned filterbank (Stage 1).

\section*{Sampling Frequency}
The sampling frequency  determines the highest representable frequency and the time resolution. For SKANN, underwater signals are band-limited to:


Thus a sampling rate of:


is chosen because:

it satisfies Nyquist for 0–32 kHz,

it is a power-of-two multiple (32k × 1.024),

it enables efficient FFT computation,

it is compatible with embedded / hydrophone ADCs.

\section*{Nyquist Frequency}
The highest representable frequency is:


\section*{For  Hz:}

Since ambient-noise PSD is defined up to 32 kHz, SKANN uses frequency mirroring for synthetic wideband signals in Stage  (Document D), but Stage 0 only processes real hydrophone data band-limited to 0–16 kHz.

FFT Size and Frequency Resolution

Given FFT size  and sampling rate , the frequency resolution is:


Typical FFT sizes in SKANN:


\section*{Examples:}


Higher  gives finer spectral detail but increases CPU/GPU load.

The FFT bin frequencies are:


Windowing and Frame Segmentation

Before computing the FFT, each frame of length  samples is multiplied by a window function  to reduce spectral leakage.

Common choices:

Hann window,

Hamming window,

Blackman window.

The windowed frame is:


Overlap and Hop Size

For continuous signals, frames overlap to ensure good time resolution:


Typical SKANN values:


This maintains smooth spectrograms and avoids information loss.

DFT Conventions and One-Sided Spectrum

Given the windowed frame , the DFT is:


For real-valued input:


Therefore, only indices  contain unique information.

Power Spectral Density (PSD) Estimation

PSD expresses mean-square pressure per unit bandwidth (Pa/Hz).

\section*{Periodogram}
The periodogram estimate is:


Where  is the window power:


The factor of 2 accounts for folding negative frequencies into positive frequencies (one-sided PSD).

Welch PSD

To obtain a smoother PSD, SKANN optionally uses Welch’s method:

Split data into overlapping frames.

Apply window to each frame.

Compute periodogram for each frame.

Average periodograms:


This reduces variance of the PSD estimate but slightly increases bandwidth.

\section*{Bandpower Computation}
Given a PSD , the mean-square pressure in band  is:


Discrete approximation:


Band level in dB:


Hermitian Symmetry and Signal Reconstruction

For reconstruction or synthesis (used in Document D), SKANN uses Hermitian symmetry to enforce real-valued time-domain signals:


\section*{With:}

This guarantees that the inverse FFT:


produces real-valued samples .

Amplitude Processing and Normalisation

Before sending any waveform into SKANN Stage 0, the amplitude must be standardised so that all datasets (synthetic and real) are comparable.

DC Offset Removal

The first step is to ensure that:


DC offset is removed by:


This prevents artificial low-frequency artefacts.

\section*{Peak Normalisation}
Peak-normalised signal:


Peak normalisation makes all recordings fit within  and is used for preliminary sanity checks.

RMS Normalisation

More physically meaningful is RMS-based scaling:


This sets the continuous-time RMS to:


Synthetic training data optionally uses target RMS values corresponding to:


\section*{Resampling Logic}
All signals must be resampled to the SKANN standard sampling rate:


\section*{Anti-Alias Filtering}
Before down-sampling by a factor :


an anti-alias filter with passband:


must be applied.

\section*{Upsampling}
Upsampling inserts zeros between samples and applies a reconstruction filter:


followed by low-pass filtering.

\section*{Polyphase Resamplers}
In practice, SKANN uses library polyphase filters (SRC / resampy) for high-quality resampling with minimal distortion.

Silence Trimming and Energy Gating

SKANN avoids frames with very low energy, since they harm SSL pretraining and classification.

A short-time energy gate is:


Frames satisfying:


(where  is a percentile threshold) are removed.

Segmentation and Padding

Segments are extracted with fixed duration:


Typical values:


Padding rules:

pad with zeros if the last segment is short,

pad symmetrically for CNNs that require even-length inputs.

Tensor Preparation for SKANN

After segmentation, each frame of length  becomes a tensor:


A batch is:


where  is batch size.

For STFT-based learners (optional):


where:



\section*{Stage 0 Preprocessing Pipeline (Summary)}
Stage 0 prepares raw or synthetic signals into a uniform digital format for ingestion by the learned filterbank (Stage 1). The steps are:

DC removal: subtract the mean to stabilise the low-frequency spectrum.

Resampling: convert all signals to  Hz with anti-alias filtering.

Amplitude normalisation: peak-normalise for safety, RMS-normalise for consistency.

Segmentation: extract segments of length  with 50\% overlap.

Windowing: apply Hann/Hamming window before FFT.

Silence gating: remove frames whose short-time energy  is below threshold.

Tensor formation: final waveform tensor:


(Optional) STFT: for spectrogram-based models:


\section*{Stage 0 Flowchart (Placeholder)}
SKANN Stage 0 preprocessing pipeline (conceptual).

Appendix A — Unified List of Symbols

The following symbols match Documents A and B to maintain consistency.

Sampling and Frames

: sampling frequency [Hz].

: sampling interval  [s].

: samples per frame.

: frame duration [s].

hop: frame advance in samples.

\section*{Frequency-Domain Quantities}
: FFT bin frequency.

: FFT bin spacing [].

: DFT of windowed frame.

: pressure PSD [Pa/Hz].

: Welch PSD estimate.

Amplitude and Energy

: discrete-time pressure.

: RMS amplitude.

: reference amplitude (Pa).

: frame energy.

\section*{Synthesis/Reconstruction}
: random phase.

IFFT: inverse discrete Fourier transform.

: amplitude spectrum.

\section*{Pandoc Compatibility Notes}
\end{document}